\section{Hermeticism}

Hermeticism is the tradition of texts attributed to Hermes Trismegistus,
centered on the Corpus Hermeticum and the Emerald Tablet. It teaches a divine
Mind as the source of all, a cosmos ordered by a word and governed by the
planetary spheres, and a human being who descends from the source and returns to
it through knowledge. Its signature doctrine is the correspondence between the
great world and the small, stated as as above, so below.

\subsection{The Creation Model}

In the beginning there is a \textbf{boundless light}, gentle and joyous, which is also \textbf{mind itself}. It is not a thing in the world. It is the ground of mind and being, the source from which everything flows. From it a part turns downward and away, condensing into a dark, churning, watery chaos. The pure and the murky are now two, the above set against the below.

Then an ordering \textbf{word} descends upon the chaos and speaks the realms apart. The fiery and light rise, the heavy and moist settle below. From the source comes a second mind, a \textbf{craftsman}, who fashions the circling heavens that bind and govern the lower world. After this the source brings forth a \textbf{luminous archetypal human} in its own image, free above the spheres and kin to the source itself. This human looks down, sees its own form reflected in matter, is drawn to it, and descends. The lower world embraces it. So the human becomes \textbf{twofold}, mortal in body yet immortal in essence. The descent comes not from sin imposed from outside but from love of one's own reflection. The way back is \textbf{knowledge}. The self climbs back up through the circling heavens, shedding at each one the passion it imposed, until it is restored to the light as pure mind.

\subsection{The One, the All, and the Mind}

Hermeticism gives several names for the highest reality.

\begin{itemize}
 \item \textbf{The One (to Hen)}, the ultimate unity from which all things come and to which they return.
 \item \textbf{The All (to Pan)}, the totality that contains everything.
 \item The divine Mind or Intellect (\textbf{Nous}), the supreme principle. The revealer names himself the Nous of the absolute sovereignty.
 \item \textbf{God as unbegotten source}, the source that itself has no source, holding all things in potential before they are made.
 \item \textbf{The Good and the Father}, the origin of life and light.
\end{itemize}

\subsection{The Poimandres Sequence}

The creation is told in the \textbf{Poimandres}, the first treatise of the \textbf{Corpus Hermeticum}, a vision granted to \textbf{Hermes Trismegistus} by the divine Mind. The sequence runs in a fixed order, each stage following from the last.

\begin{longtable}{l L{10cm} }
\caption{The Poimandres creation chain.}\\
\toprule
\textbf{Stage} & \textbf{Principle} \\
\midrule
1 & the Light, Nous, God the Father \\
2 & Darkness and watery Nature (physis) \\
3 & the Logos, the holy Word, ordering the elements \\
4 & the demiurgic Nous, the Maker, god of fire and breath \\
5 & the seven Governors of the planetary spheres \\
6 & Heimarmene, Fate, the turning of the spheres \\
7 & the Anthropos, the archetypal Human in God's image \\
8 & the fall into Nature through love of his reflection \\
9 & the twofold mortal-immortal human \\
10 & the split into male and female \\
\bottomrule
\end{longtable}

\subsection{The Seven Spheres and the Soul's Ascent}

The cosmos is built of \textbf{seven planetary spheres}, each turned by one of the \textbf{seven Governors} whose joined administration is Fate. As the soul descended into the body at birth it passed down through the spheres and took on a vice from each. At death it ascends back up and gives each vice back, gate by gate, until it stands purified.

\begin{longtable}{l L{4cm} L{6cm}}
\caption{The seven spheres and the vices the soul sheds climbing back.}\\
\toprule
\textbf{Sphere} & \textbf{Planet} & \textbf{Vice shed on ascent} \\
\midrule
1 (highest) & Saturn & increase and decrease \\
2 & Jupiter & deceit, now inactive \\
3 & Mars & desire, now powerless \\
4 & Sun & ambition of rulership \\
5 & Venus & reckless audacity \\
6 & Mercury & greed for wealth \\
7 (lowest) & Moon & ensnaring falsehood \\
\bottomrule
\end{longtable}

After the seventh gate the soul enters the \textbf{eighth nature (the Ogdoad)}, sings praise with the powers above, surrenders into them, and itself becomes a power, entering into God. This is the Hermetic deification.

\subsection{As Above, So Below}

The signature Hermetic maxim is \textbf{as above, so below}, stated in the \textbf{Emerald Tablet (Tabula Smaragdina)}. That which is below is like that which is above, and that which is above is like that which is below, to accomplish the miracles of the one thing. What holds in the heavens holds in the human and the earth.

\subsection{Macrocosm and Microcosm}

From this maxim follows the doctrine of \textbf{macrocosm and microcosm}. The human is the \textbf{microcosm}, the little world, and mirrors the cosmos, the \textbf{macrocosm}, the great world. Because the human carries the divine Mind and the archetypal Human within, it is called a great miracle (a \textbf{magnum miraculum}), a mediator joined to the gods by the divine in it and to the lower world by the body.

\subsection{The Rebirth}

The thirteenth treatise, on being born again, describes spiritual \textbf{rebirth (palingenesia)}. To be reborn one drives out the \textbf{twelve torments} bound to the body by the \textbf{ten powers} or mercies of God. As each power enters it expels the torments, until the new, divine self is born, no longer body but the immaterial intellectual self.

\subsection{An Honest Note on the Kybalion}

A widely circulated list of seven Hermetic principles, with mentalism and correspondence at its head, is \textbf{not} from the ancient Hermetica. It comes from \textbf{The Kybalion}, published in \textbf{1908} in Chicago under the pen name Three Initiates, generally identified as \textbf{William Walker Atkinson}, a New Thought writer. The book popularizes genuine Hermetic themes, above all as above so below and the primacy of Mind, but its seven-principle scheme is a modern formulation, not an ancient source text.

\subsection{Comparison to the Base Models}

Hermeticism matches the two base models of this collection closely. Like the infinite of Kabbalah it sets a source above being and lets the one descend into the many through ordered stages, and like the emptiness and luminous mind of Buddhism it makes the deepest ground a kind of Mind that the awakened self rejoins. The shared theme is clear. \textbf{The self goes out from a single source and climbs back to it.}
